\begin{tablalafrox}{Diccionario de la cuenta 5.2. Fuente: (Bases Administrativas, arts. 17.2–17.6 y Formulario T-14, p. 63); (Bases Técnicas Transversales, secciones 17.2–17.6, pp. 32–34).}{tab:T14-dic-5-2}
  {L{1.2cm} Y Y Y C{0.9cm} L{2.3cm}}
  {\cab{Código} & \cab{Paquete de trabajo} & \cab{Entregable} & \cab{Criterio de aceptación} & \cab{Resp.} & \cab{Período}}
  5.2.1 & Intercambio de lote y trazabilidad con los proveedores & Intercambio operativo con los proveedores priorizados. & El proveedor de lácteos acepta la evidencia de trazabilidad: un simulacro de retiro de uno de sus lotes en QA, con datos de prueba de recepción, lote y despacho, identifica clientes, fechas, cantidades y guías en menos de 2 horas (R18-01). & DAT & E1 · Antes del H5. \\
  5.2.2 & Integración con la pasarela de pago para el cobro con tarjeta (INT-09) & Integración en QA. & Un cobro de prueba se autoriza, se concilia con su entrega, y una falla de la pasarela no lo deja como aprobado. & DES & E1 · Antes del H4. \\
  5.2.3 & Integración con el servicio de mapas y geocodificación (INT-10) & Integración en QA. & Las rutas se generan con las direcciones georreferenciadas, y la ruta precargada funciona sin señal. & DES & E1 · Antes del H4. \\
  5.2.4 & Integración con el servicio de avisos al cliente (INT-11) & Integración en QA. & Un aviso de prueba llega por el canal configurado, y los reintentos quedan registrados. & DES & E1 · Antes del H4. \\
\end{tablalafrox}
